\chapter{Tabular Data}
This chapter focuses on tabular data and which chart idioms to use when plotting it.
As already mentioned, tabular items are identified by one or more keys and contain a vector of values specified by the attributes.
Tabular items may also have no key, which is similar to sampled values.

\section{Common Idioms for Tabular Data}
Common idioms used for tabular data include the following.
\begin{description}
    \item[Scatterplot] A scatterplot expresses quantitative attributes.
        Channels, in this case horizontal and vertical position, are linked to the attributes.
        The mark used is commonly a point.
        Scatterplots can be used to visually express and find correlation (or lack thereof), clusters, outliers, or distribution.
        A scatterplot can be even more expressive by encoding additional attributes in other channels, such as colour, size, and shape.

    \item[Bar chart] The bar chart can be used to express keyed items with commonly one quantitative attribute.
        The mark can be a line, while the channel is commonly the length of the line.
        The key can be ordered alphabetically or by attribute value.
        Bar charts can be used to compare values between keys.
    
    \item[Stacked Bar Chart] A stacked bar chart is a variation of the bar chart, where an additional key is used to express different portions of the value plotted, e.g., consumption of food (key) by country (key).
        The additional categorical channel to differentiate between keys often is colour.

    \item[Streamgraph] A streamgraph is a generalized stacked graph, where horizontal continuity is emphasized.
        One key is still channelled vertically while the other one may be continuous or a quantitative attribute.
    
    \item[Dot/Line Chart] The dot and line charts are a form of a bar chart, where the marks are connected points.
        This can be practical, when the horizontal key is quantitative or when there is a relation between nearby keys.
    
    \item[Indexed Line Chart] Indexed line charts are line charts with another key similar to stacked bar charts.
        This key is commonly encoded with colour as another line in the graph which does not have to be stacked on top.

    \item[Gantt Chart] A Gantt chart uses lines as marks which are channelled by their horizontal length to express a duration.
        The horizontal attribute commonly is time but can be any continuous attribute.
        One key can be used to differentiate between different \emph{tasks} by stacking them vertically.

    \item[Slopegraph] A slopegraph uses two quantitative attributes, but only samples two points on the horizontal (often time) attribute.
        Each left and right sample is marked using points which are connected.
        The vertical attribute is associated with the key and shown on the left and right size of the two points.
        The change of the vertical attribute (positive or negative) is often also channelled using colour or/and size.

    \item[Parallel Coordinates] parallel coordinate charts are slopegraphs with more than two points.

    \item[Heatmap] For tables with two keys (can be continuous, e.g., space) and one quantitative attribute, heatmaps display values in a 2D matrix which is indexed by the two keys.
        Colour is used to channel the values.
        Heatmaps can be used to find clusters and outliers.

    \item[Cluster Heatmap] A cluster heatmap is a heatmap, where the keys are ordered by identified clusters.

    \item[Radial Bar Chart, Star Plot] Radial bar charts and star plots are radial bar charts.

    \item[Radar Plot] A radar plot is a radial line chart.
        Should be avoided unless the categorical or quantitative key is radial.

    \item[Pie Chart, Coxcomb Chart] Pie charts and coxcomb charts divide the area of a circle into different regions based on a categorical key.
        The size of each area is linked to the quantitative attribute to map.
        Similar to stacked bar charts, these charts can also be stacked to channel even another key.
        Small differences can be recognized easier when using coxcomb charts rather than pie charts.

    \item[Normalized Stacked Bar Chart] this chart is a stacked bar chart, where the ratio of the second key is displayed by normalizing each stacked value to full vertical height.
        Each line is equivalent to one entire pie chart.
\end{description}
